\chapter{System Reproduction Guide and Runbook}
\label{app:runbook}

This appendix provides a step-by-step operational runbook to reproduce all software builds, automated test suites, Mininet network emulations, and agentic AI features documented in this report.

\section{System Prerequisites}
The software stack requires a Linux environment with root network namespace permissions:
\begin{itemize}
  \item \textbf{Operating System}: Ubuntu 22.04 LTS or Arch Linux (Kernel $\ge 5.15$).
  \item \textbf{Python}: Python 3.9 or 3.10 with virtual environment tools.
  \item \textbf{Node.js \& npm}: Node.js $\ge 18.0$ (LTS) and npm $\ge 9.0$.
  \item \textbf{Networking Packages}: \texttt{mininet}, \texttt{openvswitch-switch}, \texttt{iproute2}, \texttt{iperf3}.
  \item \textbf{LaTeX Tooling}: \texttt{texlive-latex-base}, \texttt{texlive-latex-extra}, \texttt{texlive-fonts-recommended}.
\end{itemize}

\section{Environment Installation}
Clone the repository and initialize the Python virtual environment:
\begin{lstlisting}[language=bash]
# 1. Clone repository and enter directory
cd /opt/sdn-controller  # or path to project repository

# 2. Activate virtual environment
source .venv/bin/activate

# 3. Install core Python dependencies
pip install -r requirements.txt

# 4. Install frontend npm dependencies
npm --prefix future/dashboard/react-app install
\end{lstlisting}

\section{Running Automated Test Suites}
The project enforces strict hermetic testing:

\subsection{Unit and API Test Suite}
Execute the complete pytest suite:
\begin{lstlisting}[language=bash]
PYTHONPATH=".:future" pytest tests/ -v
\end{lstlisting}
\textbf{Expected Result}: \texttt{17 passed} (100\% pass rate in $\sim$75 seconds).

\subsection{Frontend TypeScript and Vite Build}
Verify clean compilation of the React dashboard:
\begin{lstlisting}[language=bash]
npm --prefix future/dashboard/react-app run build
\end{lstlisting}
\textbf{Expected Result}: \texttt{tsc -b \&\& vite build} exits with code 0 (zero TypeScript errors).

\section{Launching the Live System}
The live interactive system is executed across three terminal daemons:

\subsection{Daemon 1: Ryu OpenFlow 1.3 Controller}
\begin{lstlisting}[language=bash]
PYTHONPATH=".:future" ryu-manager \
  --ofp-listen-host 127.0.0.1 \
  --ofp-tcp-listen-port 6653 \
  controller.sdn_controller
\end{lstlisting}

\subsection{Daemon 2: FastAPI Backend Service}
\begin{lstlisting}[language=bash]
PYTHONPATH=".:future" python -m uvicorn future.api.main:app \
  --host 127.0.0.1 --port 8000
\end{lstlisting}

\subsection{Daemon 3: React Development Server}
\begin{lstlisting}[language=bash]
npm --prefix future/dashboard/react-app run dev
\end{lstlisting}
Open a web browser at \url{http://localhost:5173} to access the live dashboard.

\section{Testing Autonomous Congestion and Remediation}
To test the LangGraph agentic loop:
\begin{enumerate}
  \item On the top navigation bar, click the \textbf{Agentic AI} modal icon.
  \item In the chat modal, click the quick action pill: \texttt{[Create Congestion (14M)]}.
  \item Observe link $s_1 \leftrightarrow s_2$ surging to 88\% on the canvas and strip chart.
  \item Type or click \texttt{[Fix Congestion]}.
  \item Verify that the 5-stage LangGraph execution pipeline transitions to \texttt{COMPLETED}, installs Make-Before-Break flow rules, and validates telemetry recovery to 1.0\%.
\end{enumerate}

\section{Compiling this Technical Report}
To compile this 30--35 page LaTeX technical report into high-resolution PDF format:
\begin{lstlisting}[language=bash]
cd report/
pdflatex -interaction=nonstopmode main.tex
pdflatex -interaction=nonstopmode main.tex
\end{lstlisting}
Two passes are executed to ensure complete resolution of cross-references, table of contents, list of figures, and list of tables. The output document is generated as \texttt{report/main.pdf}.
